\section{Split verification sheet}
\label{sec:appendix-verification-sheet}

% Source: runs/20260814_figure_verification_sheet_02/. The figure is a
% photographic contact sheet and is PNG by decision, not by omission: it embeds
% 24 real thumbnails, and vector output would carry the same rasters at a larger
% file size. See figures/README.md.
%
% RE-RENDERED 2026-08-14 for the page. It was 12 rows at 13.5 in wide, which
% reduced to about 46% at text width and put its smallest text near 3.7 pt.
% It is now 6 rows at 7.0 in with all text at 8.0 pt, sized for the ~7.16 in
% text width of a two-column IEEE page. Under THIS class the page gives 6.27 in,
% so the text lands at about 7.2 pt; under the venue template it is roughly 1:1.
% Becomes figure* there -- it already is one here, which costs nothing in a
% single-column class.

Figure~\ref{fig:verification-sheet} is the visual counterpart to the
contamination figures quoted in Sections~\ref{sec:duplicate-contamination}
and~\ref{sec:correction-limits}. Those sections report that no test--train pair
qualifies as a near-duplicate at any threshold tested, and that the closest such
pair clears the loosest threshold by 3\%. A reader has no way to judge from a
count alone whether a scoring function that returns zero is measuring the right
thing, so the sheet shows the pairs the scorer ranked highest and lets the claim
be checked by eye rather than taken on trust.

Two properties of the sheet make it evidence rather than illustration. It scores
every cross-split pair with \emph{no} epsilon cutoff, so the right-hand block
shows the six closest test--train pairs even though none of them qualifies as a
duplicate --- non-qualifying pairs are the point, not an omission. And the
per-pair distances are printed beside each thumbnail, so the 0.0515 margin
behind the negative result is visible as a number and as two photographs at the
same time.

\begin{figure*}[p]
  \centering
  \includegraphics[width=\textwidth,height=0.86\textheight,keepaspectratio]{split_verification_sheet.png}
  \caption{Split verification sheet. Left: the six most similar cross-split
  pairs across all three relationships. Right: the six closest test--train
  pairs, none of which qualifies as a near-duplicate at any threshold tested.
  Every cross-split pair is scored with no epsilon cutoff, so the right-hand
  block exists precisely because its pairs do not qualify. Of the 4{,}860
  test--train pairs sharing a class inventory, 4{,}837 carry more than two
  annotated boxes and are ranked; 23 low-information pairs are excluded from
  that ranking, of the 95 excluded across all relationships. The closest
  test--train pair scores 0.0515 against a loosest threshold of 0.05.}
  \label{fig:verification-sheet}
\end{figure*}
